% page 39 du fac-simile (image p-038 du scan)
% bloc 162.6 x 246.3 mm ; marge gauche 39.4 mm
\pgc{17.780mm}{0.000mm}{
\l{}
\l{\cel{50}31}
\l{}
\l{}
\l{\sou{antifermento}. (\sou{kemio, biol}.)Korpo qua su opozas a la efiko di}
\l{\cel{3}fermento. - DEFIS.}
\l{}
\l{\sou{antifono}. (\sou{liturgio katol}.) Versaro quan onu kantas, tote o}
\l{\cel{3}parte, ante psalmo o kantiko, e quan onu repetas tote,pose.}
\l{\cel{3}- DEFIRS.}
\l{}
\l{\sou{antifrazo}. Frazeto uzata - eufemisme od ironie - por expresar}
\l{\cel{3}la kontreajo di la senco vera. - DEFIRS.}
\l{}
\l{\sou{antiklina}.\cel{2}(\sou{geol.}) Dicesas pri la plisuro di strato en qua}
\l{\cel{3}la parto konkava regardas adinfre.}
\l{}
\l{\sou{antilogio}. Des-kohero di la idei, di la dico-maniero.-DEFIS.}
\l{}
\l{\sou{antilopo}.\cel{2}(\sou{zool}.) Ruminero di formo gracila, gracoza, kun}
\l{\cel{3}korni ne-kaduka,\cel{1}\sou{kava,}\cel{1}cirkondanta, nukleo ostoza, bloka.}
\l{\cel{3}- L.\cel{1}\sou{antilope}.- DEFIRS.}
\l{}
\l{\sou{antimonio}.\cel{2}(\sou{kemio)}.\cel{2}Korpo simpla, metala, qua uzesas por}
\l{\cel{2}kompozar emetiko, e qua pulverigite uzesas nomo di \textquotedbl{}kohi\textquotedbl{}.}
\l{\cel{3}- DEFIRS.}
\l{}
\l{\sou{antinomio}. Des-kohero, reala o semblanta, di du legi naturala}
\l{\cel{3}(+leyi) o yurala. - DEFIRS.}
\l{}
\l{\sou{antipatiar}.\cel{2}(\sou{trans.}) Instinte eskartar su de (ulo o de ulu).}
\l{\cel{3}- DEFIRS.}
\l{}
\l{\sou{antipirino}.- Medikamento obtenita de un ek la kompozanti di}
\l{\cel{3}karbono, e qua kombatas febro, nevralgii. - DEFIRS.}
\l{}
\l{\sou{antipodo}. I. Persono qua okupas sur la Terglobo loko diametre}
\l{\cel{3}opozata. - II. Sur la Terglobo, loko diametre opozata ad}
\l{\cel{3}altra. - DEFIRS.}
\l{}
\l{\sou{antiqua}.\cel{2}Qua apartenas a la vivo-maniero, a la kustumi di epoki}
\l{\cel{3}tre antea (partikulare, qua apartenas a la civilizeso-sistemo}
\l{\cel{3}di la Greki e di la Romani lor la quaresma e lor la unesma}
\l{\cel{3}yarcenti rispektive ante nia ero). - DEFIRS.}
\l{}
\l{\sou{antisepsiar}. (\sou{trans.}) Uzar produkturo apta kombatar la putro}
\l{\cel{3}morbala. - DEFIRS.}
\l{}
\l{\sou{antitezo}. Kontrasto qua produktesas da la apudigo di du idei,}
\l{\cel{3}di du expresuri, qui inter-opozesas. - DEFIRS.}
\l{}
\l{\sou{antitoxino}. Substanco apta nihiligar la efiko di toxino. - DEF.}
\l{}
\l{\sou{antologio}. Kolektajo de mikra versari. - DEFIRS.}
\l{}
\l{\sou{antonimo}. Vorto qua, konsiderate relate ad altra vorto, havas}
\l{\cel{3}senco absolute opozata. - DEFIRS.}
}
